% ---------------------------------------------------------------------------
\chapter{Prospetto dei quaderni}
\label{app:prospetto}

La tabella elenca i ventitr\'e file nell'ordine in cui compaiono nel volume.
La colonna \emph{Espr.} d\`a il numero di espressioni recuperate, la colonna
\emph{Verif.} quante di esse hanno superato il controllo numerico descritto
nel capitolo~\ref{cap:verifica}. Si ricordi che un valore basso in
quest'ultima colonna dice soprattutto che il quaderno contiene equazioni da
risolvere o calcoli di lavoro, non identit\`a da controllare.

\begin{longtable}{@{}l l r r r@{}}
\toprule
\textbf{File} & \textbf{Data di salvataggio} & \textbf{Espr.} & \textbf{Verif.} & \textbf{Pag.} \\
\midrule
\endfirsthead
\toprule
\textbf{File} & \textbf{Data di salvataggio} & \textbf{Espr.} & \textbf{Verif.} & \textbf{Pag.} \\
\midrule
\endhead
\texttt{teoria.dfw} & 06 May 2009 & 27 & 0 & \pageref{cap:teoria} \\
\texttt{teoria2.dfw} & 02 Nov 2009 & 31 & 0 & \pageref{cap:teoria2} \\
\texttt{limit.dfw} & 19 Dec 2009 & 1 & 0 & \pageref{cap:limit} \\
\texttt{e.dfw} & 27 Aug 2010 & 93 & 15 & \pageref{cap:e} \\
\texttt{pigreco.dfw} & 27 Aug 2010 & 26 & 2 & \pageref{cap:pigreco} \\
\texttt{sin-cos.dfw} & 04 Jan 2010 & 20 & 9 & \pageref{cap:sin-cos} \\
\texttt{j.dfw} & 07 Jul 2009 & 7 & 3 & \pageref{cap:j} \\
\texttt{atan.dfw} & 28 Jan 2010 & 14 & 8 & \pageref{cap:atan} \\
\texttt{equazioni.dfw} & 27 Jan 2010 & 44 & 0 & \pageref{cap:equazioni} \\
\texttt{mate.dfw} & 18 Aug 2010 & 133 & 35 & \pageref{cap:mate} \\
\texttt{mate copia.dfw} & 09 Feb 2010 & 127 & 34 & \pageref{cap:mate-copia} \\
\texttt{cosinc.dfw} & 14 Nov 2010 & 29 & 3 & \pageref{cap:cosinc} \\
\texttt{cost\textasciicircum{}2.dfw} & 24 Jun 2010 & 22 & 0 & \pageref{cap:cost2} \\
\texttt{gauss.dfw} & 15 Sep 2010 & 14 & 3 & \pageref{cap:gauss} \\
\texttt{delta.dfw} & 13 Jul 2010 & 55 & 7 & \pageref{cap:delta} \\
\texttt{Trasformata T.dfw} & 17 Jan 2011 & 14 & 13 & \pageref{cap:trasformata-t} \\
\texttt{fx\textasciicircum{}gx.dfw} & 17 Nov 2010 & 10 & 0 & \pageref{cap:fxgx} \\
\texttt{factor.dfw} & 18 Aug 2010 & 36 & 13 & \pageref{cap:factor} \\
\texttt{fattorazionale.dfw} & 07 Jan 2011 & 14 & 0 & \pageref{cap:fattorazionale} \\
\texttt{numeri primi.dfw} & 27 Aug 2010 & 1 & 0 & \pageref{cap:numeri-primi} \\
\texttt{u(0)-sign(0).dfw} & 04 Jan 2011 & 12 & 0 & \pageref{cap:u0-sign0} \\
\texttt{Sudoku.dfw} & 26 Apr 2009 & 42 & 0 & \pageref{cap:sudoku} \\
\texttt{notch.dfw} & 21 Jan 2010 & 20 & 0 & \pageref{cap:notch} \\
\midrule
\textbf{Totale} & & \textbf{792} & \textbf{145} & \\
\bottomrule
\end{longtable}


% ---------------------------------------------------------------------------
\chapter{I sorgenti del recupero}
\label{app:sorgenti}

Quello che segue \`e il codice che ha prodotto tutto il resto del volume.
\`E scritto in Python 3 e non dipende da nulla al di fuori di
\texttt{numpy}, \texttt{matplotlib} e \texttt{mpmath}.

\section{Estrazione dal formato binario}
\lstinputlisting[language=Python]{../tools/dfw_extract.py}

\section{Traduzione in \LaTeX}
\lstinputlisting[language=Python]{../tools/derive2tex.py}

\section{Verifica numerica}
\lstinputlisting[language=Python]{../tools/verify.py}

\section{Generazione delle figure}
\lstinputlisting[language=Python]{../tools/make_figures.py}

\section{Composizione dei capitoli}
\lstinputlisting[language=Python]{../tools/build_tex.py}
